\chapter{Medicines for pain}
\index{Medicines for pain}

\section*{Definition and Overview}
Medicines for pain, or analgesics, are a diverse class of pharmaceutical agents designed to alleviate pain by acting on various components of the peripheral and central nervous systems. Pain management is a fundamental aspect of clinical practice, categorised by duration (acute versus chronic) and underlying pathophysiology (nociceptive, neuropathic, or nociplastic). The pharmacological options for pain relief span several classes, including simple non-opioid analgesics (such as paracetamol), non-steroidal anti-inflammatory drugs (NSAIDs, including ibuprofen and celecoxib), opioid receptor agonists (such as morphine, oxycodone, and fentanyl), and adjuvant medications (such as anticonvulsants like pregabalin, and antidepressants like amitriptyline). Selecting the appropriate analgesic regimen requires a comprehensive understanding of each drug's mechanism of action, pharmacokinetics, and safety profile to tailor treatment to the individual patient, balancing efficacy with the risk of adverse effects and substance misuse.

\section*{Epidemiology}
Pain is one of the most common reasons for seeking medical care worldwide and within the affected region. Chronic pain alone affects approximately 1 in 5 people aged 45 and over, with the prevalence increasing with age. Musculoskeletal disorders, such as osteoarthritis and chronic lower back pain, are the leading causes of chronic pain. In contrast, acute pain is a near-universal human experience associated with trauma, surgery, or acute medical illness. The epidemiology of pain management has shifted significantly over the past two decades due to the opioid crisis, leading to stricter clinical guidelines and regulatory measures globally and in many countries. Despite these restrictions, prescription and over-the-counter analgesics remain among the most commonly dispensed medicines. The use of adjuvant analgesics (such as gabapentinoids) has risen concurrently as clinicians seek non-opioid alternatives for neuropathic and chronic pain conditions.

\section*{Aetiology and Pathophysiology}
Pain perception involves complex neurophysiological processes, and different pain medicines are selected based on the specific physiological mechanism they target within the pain pathways.
\begin{itemize}
    \item \textbf{Nociceptive pain pathway:} In response to tissue damage, peripheral nociceptors are activated by chemical mediators (such as prostaglandins, bradykinin, and histamine). These signals travel along A-delta and C fibres to the dorsal horn of the spinal cord, where they synapse and ascend via the spinothalamic tract to the thalamus and cerebral cortex.
    \item \textbf{Mechanism of paracetamol:} Paracetamol exhibits analgesic and antipyretic properties, but minimal anti-inflammatory activity. Its precise mechanism of action remains complex, involving the inhibition of cyclooxygenase (COX) enzymes in the central nervous system, modulation of the endogenous cannabinoid system, and activation of descending serotonergic inhibitory pathways.
    \item \textbf{Mechanism of NSAIDs:} Non-steroidal anti-inflammatory drugs inhibit cyclooxygenase enzymes (COX-1 and/or COX-2). COX-2 inhibition reduces the synthesis of pro-inflammatory prostaglandins (such as PGE2) at the site of tissue injury, decreasing nociceptor sensitisation. COX-1 inhibition, however, reduces cytoprotective prostaglandins in the gastric mucosa and impairs platelet aggregation.
    \item \textbf{Mechanism of opioid analgesics:} Opioids act as agonists at G-protein coupled opioid receptors (primarily mu-opioid receptors) in the brain, spinal cord, and gastrointestinal tract. Activation inhibits presynaptic neurotransmitter release (glutamate, substance P) and hyperpolarises postsynaptic neurons via potassium channel activation, effectively blocking transmission of pain signals and activating descending inhibitory pathways.
    \item \textbf{Mechanism of adjuvant analgesics:} Anticonvulsants (such as gabapentin and pregabalin) bind to the alpha-2-delta subunit of voltage-gated calcium channels in the central nervous system, reducing calcium influx and subsequent release of excitatory neurotransmitters in neuropathic pain. Tricyclic antidepressants (such as amitriptyline) and serotonin-noradrenaline reuptake inhibitors (SNRIs, such as duloxetine) enhance descending pain inhibitory pathways by preventing the reuptake of serotonin and noradrenaline.
\end{itemize}

\section*{Clinical Features}
The clinical presentation of pain and the therapeutic effects of pain medicines depend on the type of pain (nociceptive versus neuropathic) and the pharmacological characteristics of the prescribed agent.
\begin{itemize}
    \item \textbf{Nociceptive pain characteristics:} Typically described as dull, aching, throbbing, or sharp, and is well-localised to the area of tissue damage (for example, postsurgical pain, osteoarthritis, or acute injury).
    \item \textbf{Neuropathic pain characteristics:} Characterised by burning, shooting, stabbing, electric-shock-like sensations, often accompanied by allodynia (pain from a non-painful stimulus) and hyperalgesia (increased sensitivity to pain).
    \item \textbf{Therapeutic response to non-opioids:} Reduction in baseline pain intensity and inflammation, often sufficient for mild-to-moderate nociceptive pain and acute inflammatory states.
    \item \textbf{Therapeutic response to opioids:} Rapid and profound analgesia, along with an altered emotional response to pain (reduced distress), typically reserved for severe acute pain, post-operative recovery, or cancer pain.
    \item \textbf{Therapeutic response to adjuvants:} Gradual reduction in neuropathic pain severity, shooting pain frequency, and paraesthesias, often requiring several weeks of consistent dosing to achieve full efficacy.
    \item \textbf{Common side effects:} Non-opioids can cause dyspepsia or renal impairment; opioids frequently cause constipation, nausea, somnolence, respiratory depression, and miosis; adjuvants commonly cause dizziness, peripheral oedema, and cognitive clouding.
\end{itemize}

\section*{Differential Diagnosis}
Prior to initiating analgesic therapy, a precise diagnosis of the underlying cause of pain is essential, as managing pain as a symptom must not delay or mask the treatment of a serious, treatable primary condition.
\begin{itemize}
    \item \textbf{Nociceptive versus neuropathic pain:} Discerning the primary pain mechanism is critical, as nociceptive pain responds well to simple analgesics and NSAIDs, whereas neuropathic pain requires adjuvants like anticonvulsants or antidepressants.
    \item \textbf{Referred pain versus localised pain:} Pain in a specific dermatome or body region must be traced to its origin (such as myocardial infarction presenting as jaw or arm pain, or radiculopathy presenting as leg pain) rather than treating only the symptomatic site.
    \item \textbf{Chronic primary pain versus chronic secondary pain:} Chronic primary pain (such as fibromyalgia, where pain is the disease) must be differentiated from secondary pain (such as pain secondary to active rheumatoid arthritis or cancer), which requires targeting the underlying inflammatory or neoplastic pathology.
    \item \textbf{Opioid-induced hyperalgesia versus tolerance:} In patients on long-term opioid therapy, a paradoxically increased sensitivity to pain (hyperalgesia) must be distinguished from pharmacological tolerance, as hyperalgesia requires opioid tapering rather than dose escalation.
    \item \textbf{Psychosocial contributors to pain:} Conditions such as major depressive disorder, somatisation, or high levels of pain catastrophising must be identified and managed, as they profoundly influence pain perception and response to pharmacotherapy.
\end{itemize}

\section*{When to Seek Medical Care}
Patients experiencing new, worsening, or unexplained pain should seek a clinical evaluation to determine the underlying cause and establish an appropriate management plan. Immediate emergency services should be contacted for severe, life-threatening symptoms associated with acute pain or analgesic toxicity.

\textbf{Contact your local emergency services for:}
\begin{itemize}
    \item Sudden, crushing, or squeezing chest pain that may radiate to the jaw, neck, back, or arm
    \item Severe, sudden-onset headache (``thunderclap'' headache) or pain accompanied by neck stiffness and high fever
    \item Sudden, severe, and unexplained abdominal pain, especially if accompanied by persistent vomiting, abdominal rigidity, or high fever
    \item Signs of opioid toxicity or overdose, including severe respiratory depression (slowed breathing), pinpoint pupils, and inability to be aroused
    \item Sudden weakness, numbness, or loss of bowel or bladder control associated with back pain
\end{itemize}

\section*{Diagnosis and Investigations}
A comprehensive diagnostic workup is essential to classify the pain type, identify the underlying pathology, and rule out contraindications to specific analgesic medications.
\begin{itemize}
    \item \textbf{Detailed pain history and assessment scales:} Evaluating pain location, quality, duration, exacerbating/relieving factors, and using validated scales to measure severity and functional impact.
    \item \textbf{Diagnostic imaging:} X-rays, ultrasound, CT, or MRI scans are performed to identify structural abnormalities, such as fractures, joint degeneration, disk herniations, or visceral pathology.
    \item \textbf{Electrophysiological studies:} Nerve conduction studies and electromyography (NCS/EMG) are indicated to confirm and localise suspected peripheral nerve damage or radiculopathy.
    \item \textbf{Renal and hepatic function tests:} Baseline urea, creatinine, and liver transaminases must be assessed, as renal impairment restricts NSAID use and hepatic impairment requires caution or dose reductions for paracetamol and certain opioids.
    \item \textbf{Cardiovascular risk assessment:} Evaluating for hypertension, ischaemic heart disease, or heart failure prior to prescribing NSAIDs, particularly selective COX-2 inhibitors, which increase the risk of cardiovascular events.
\end{itemize}

\section*{Management}
Analgesic management should follow a multimodal and multidisciplinary approach, combining pharmacological therapies with non-pharmacological, physical, and psychological strategies tailored to the individual patient's pain profile.
\begin{itemize}
    \item \textbf{Stepwise pharmacological escalation:} Utilising a structured approach starting with non-opioid medications (paracetamol, NSAIDs) for mild pain, and adding weak or strong opioids only for severe acute pain or cancer-related pain, while monitoring for efficacy and side effects.
    \item \textbf{Adjuvant medication selection:} Initiating low-dose tricyclic antidepressants, SNRIs, or gabapentinoids for patients with confirmed neuropathic pain, gradually titrating the dose over several weeks.
    \item \textbf{Non-pharmacological and physical therapies:} Incorporating physiotherapy, structured exercise programmes, heat/cold applications, and occupational therapy to improve physical function and support rehabilitation.
    \item \textbf{Psychological and behavioural interventions:} Utilising cognitive behavioural therapy (CBT), mindfulness, and pain education to help patients develop coping strategies, reduce anxiety, and manage the impact of chronic pain on their daily lives.
    \item \textbf{Opioid stewardship and safety monitoring:} Implementing strict monitoring when prescribing opioids, including prescribing short courses for acute pain, establishing clear treatment agreements, utilising real-time prescription monitoring systems, and co-prescribing naloxone for patients at risk of overdose.
\end{itemize}

\section*{Complications}
\begin{itemize}
    \item \textbf{Opioid use disorder and dependency:} The development of physical dependence, tolerance, and addiction (substance use disorder) following prolonged opioid therapy.
    \item \textbf{Gastrointestinal ulceration and haemorrhage:} Gastric mucosal damage, perforation, or serious bleeding caused by chronic or high-dose NSAID use.
    \item \textbf{Analgesic-induced nephropathy and acute kidney injury:} Declines in renal function, fluid retention, and acute kidney injury, particularly when NSAIDs are used in patients with pre-existing renal impairment or dehydration.
    \item \textbf{Medication overuse headache:} A chronic headache condition triggered by the frequent or daily use of acute headache medications (such as paracetamol, triptans, or opioids) for more than 10 to 15 days per month.
    \item \textbf{Hepatotoxicity:} Liver damage or acute liver failure resulting from paracetamol overdose, either accidental (through taking multiple products containing paracetamol) or intentional.
\end{itemize}

\section*{Prognosis and Outcomes}
The prognosis and clinical outcomes for patients requiring pain medicines depend heavily on the underlying cause of the pain, its duration, and the management strategy employed. Acute pain typically has an excellent prognosis, resolving completely as the underlying tissue injury heals, with analgesic requirements decreasing rapidly. In contrast, the prognosis for chronic pain is more complex. While complete pain elimination is rarely achieved in chronic conditions, a multimodal treatment plan can successfully reduce pain intensity, improve physical function, and enhance quality of life. Early intervention and the avoidance of long-term, high-dose opioid therapy are associated with better long-term functional outcomes and lower risks of dependency or adverse drug events. Interdisciplinary pain management programmes that focus on self-management and functional goals yield the most sustainable outcomes for chronic pain syndromes.

\section*{Prevention and Self-Care}
\begin{itemize}
    \item \textbf{Correct dosing and avoiding double-dosing:} Always check the active ingredients of all medicines you take to avoid accidentally double-dosing on the same class, such as taking two different paracetamol-containing products.
    \item \textbf{Use the lowest effective dose for the shortest time:} Minimise the risk of side effects by taking pain relievers, particularly NSAIDs and opioids, only when necessary and at the lowest dose that provides relief.
    \item \textbf{Take NSAIDs with food:} Reduce the risk of stomach irritation and indigestion by taking non-steroidal anti-inflammatory drugs with food or a glass of milk.
    \item \textbf{Prevent opioid-induced constipation:} Increase dietary fibre, maintain adequate hydration, and engage in regular light physical activity, or use a prescribed laxative if taking opioid medicines.
    \item \textbf{Safely discard unused medicines:} Return any leftover prescription pain medicines, especially opioids, to a local pharmacy for safe disposal to prevent accidental poisoning or misuse by others.
\end{itemize}

\section*{Sources and Further Reading}
The following reputable organisations provide further information on pain medicines and their safe administration:
\begin{itemize}
    \item Therapeutic Goods Administration. \textit{Information on prescription pain medicines and safety regulations in Australia.} https://www.tga.gov.au/
    \item Painaustralia. \textit{National advocacy and support organisation for people living with pain.} https://www.painaustralia.org.au/
    \item Healthdirect Australia. \textit{Guide to choosing and using pain relief medicines safely.} https://www.healthdirect.gov.au/
    \item Faculty of Pain Medicine, ANZCA. \textit{Clinical guidelines and resources on professional pain management.} https://www.fpm.anzca.edu.au/
    \item NHS. \textit{Guide to common painkillers and managing chronic pain.} https://www.nhs.uk/
\end{itemize}
